\chapter{具身外部记忆与社会性外部记忆}
\chaptervoice{这一章，我们围绕“具身外部记忆与社会性外部记忆”做一件克制的事：把直觉压缩为状态、关系、时间尺度和可检验后果。它在全书中的任务，是说明固定功能拓扑中的持续主体。}

\section*{理论来路与依赖接口}
本章的直接思想来源是原文关于 CCM、调度、边界、卸载、单主体与 L0—L6 多层循环的 AgentOS 讨论。我们采用原文较晚形成的定义来整理较早讨论，不把对话中的试探性比喻与最终模型并列成互相竞争的“定律”。已有理论锚点主要是持续学习、工作空间、观测器—控制器与混杂系统（Kalman、混杂系统理论与 Cowan）。

本章使用上一章“内部记忆、外部记忆与 Me”已经得到的对象与边界；本章产出将供下一章“连续动力学与离散治理”使用。若后文术语偶尔出现，只作为路线提示，不参与本章推导。

\section*{从哪里出发}
我们已经拥有上一章给出的概念，现在只向前走一步。我会先把对象放进闭环，再讨论它的数学形式，最后检查工程边界。读者不必预先相信本章术语；只需暂时接受一个方法论约定：智能结构的意义来自它在现实反馈中的作用，而不是来自名称本身。

令系统在观察窗口 $[t,t+T]$ 内的状态轨迹为 $x_{[t,t+T]}$，环境扰动为 $d$，行动为 $u$。本章讨论的对象若确实有效，就应改变条件分布 $p(x_{t+T}\mid x_t,u,d)$，或改变达到同一结果所需的代价。这个起点把讨论锚定在持续学习、工作空间、观测器—控制器与混杂系统上。
\section{修改环境就是写外部记忆}
\begin{narration}
现在把镜头移到“修改环境就是写外部记忆”。我们只沿本章已经取得的概念向前一步：先确定它在闭环里的角色，再检查它的尺度与失败边界。
\end{narration}

本节把“修改环境就是写外部记忆”落实为跨主体 CSO 流通过角色、协议、外部记忆和制度沉降形成的慢尺度协调结构。这个定义不是说“任何相似现象都算”，而是要求我们同时给出对象域、允许的干预以及观察窗口。对于主题“具身外部记忆与社会性外部记忆”而言，它承担的是固定功能拓扑中的持续主体中的一个可辨识环节。

把这一关系写成最小数学骨架，可得

\begin{equation}\mathcal O=(\{A_i\},\mathcal P,\mathcal M,\mathcal N).\end{equation}

式中的符号只承担结构角色，具体参数仍需由任务和身体数据辨识。我们首先检查可观测量，再检查控制量，最后检查比它更慢的约束。这样的顺序来自持续学习、工作空间、观测器—控制器与混杂系统，但本书对修改环境就是写外部记忆的命名和组合仍是需要实验验证的建模提案。

这一小节的反例同样重要：把组织类比成单一大脑会虚构不存在的中央观察器。因此评价它不能只看一次任务结果，而要比较干预前后的轨迹、恢复时间、维护成本和风险。如果改变该机制而所有允许观察下的分布都不变，那么它至少在当前实验族中是冗余描述。

\begin{thought}
对“修改环境就是写外部记忆”做一次消融：冻结它、删除它，或只允许它以相邻时间尺度交换残差。哪些现象立即消失，哪些只在长期出现？若答案无法区分三种干预，我们还没有找到正确的状态变量。
\end{thought}

最后，把结论交还现实：本节只建立功能层关系，不由此推出特定脑区实现、主观意识或宇宙目的。可测状态、干预后果和明确失败条件，才是从原文直觉走向可检验理论的桥。

\section{工具是稳定的动作结构}
\begin{narration}
现在把镜头移到“工具是稳定的动作结构”。我们只沿本章已经取得的概念向前一步：先确定它在闭环里的角色，再检查它的尺度与失败边界。
\end{narration}

本节把“工具是稳定的动作结构”落实为跨主体 CSO 流通过角色、协议、外部记忆和制度沉降形成的慢尺度协调结构。这个定义不是说“任何相似现象都算”，而是要求我们同时给出对象域、允许的干预以及观察窗口。对于主题“具身外部记忆与社会性外部记忆”而言，它承担的是固定功能拓扑中的持续主体中的一个可辨识环节。

把这一关系写成最小数学骨架，可得

\begin{equation}\mathcal O=(\{A_i\},\mathcal P,\mathcal M,\mathcal N).\end{equation}

式中的符号只承担结构角色，具体参数仍需由任务和身体数据辨识。我们首先检查可观测量，再检查控制量，最后检查比它更慢的约束。这样的顺序来自持续学习、工作空间、观测器—控制器与混杂系统，但本书对工具是稳定的动作结构的命名和组合仍是需要实验验证的建模提案。

这一小节的反例同样重要：把组织类比成单一大脑会虚构不存在的中央观察器。因此评价它不能只看一次任务结果，而要比较干预前后的轨迹、恢复时间、维护成本和风险。如果改变该机制而所有允许观察下的分布都不变，那么它至少在当前实验族中是冗余描述。

\begin{thought}
对“工具是稳定的动作结构”做一次消融：冻结它、删除它，或只允许它以相邻时间尺度交换残差。哪些现象立即消失，哪些只在长期出现？若答案无法区分三种干预，我们还没有找到正确的状态变量。
\end{thought}

最后，把结论交还现实：本节只建立功能层关系，不由此推出特定脑区实现、主观意识或宇宙目的。可测状态、干预后果和明确失败条件，才是从原文直觉走向可检验理论的桥。

\section{其他 Agent 作为认知资源}
\begin{narration}
现在把镜头移到“其他 Agent 作为认知资源”。我们只沿本章已经取得的概念向前一步：先确定它在闭环里的角色，再检查它的尺度与失败边界。
\end{narration}

本节把“其他 Agent 作为认知资源”落实为跨主体 CSO 流通过角色、协议、外部记忆和制度沉降形成的慢尺度协调结构。这个定义不是说“任何相似现象都算”，而是要求我们同时给出对象域、允许的干预以及观察窗口。对于主题“具身外部记忆与社会性外部记忆”而言，它承担的是固定功能拓扑中的持续主体中的一个可辨识环节。

把这一关系写成最小数学骨架，可得

\begin{equation}\mathcal O=(\{A_i\},\mathcal P,\mathcal M,\mathcal N).\end{equation}

式中的符号只承担结构角色，具体参数仍需由任务和身体数据辨识。我们首先检查可观测量，再检查控制量，最后检查比它更慢的约束。这样的顺序来自持续学习、工作空间、观测器—控制器与混杂系统，但本书对其他 Agent 作为认知资源的命名和组合仍是需要实验验证的建模提案。

这一小节的反例同样重要：把组织类比成单一大脑会虚构不存在的中央观察器。因此评价它不能只看一次任务结果，而要比较干预前后的轨迹、恢复时间、维护成本和风险。如果改变该机制而所有允许观察下的分布都不变，那么它至少在当前实验族中是冗余描述。

\begin{thought}
对“其他 Agent 作为认知资源”做一次消融：冻结它、删除它，或只允许它以相邻时间尺度交换残差。哪些现象立即消失，哪些只在长期出现？若答案无法区分三种干预，我们还没有找到正确的状态变量。
\end{thought}

最后，把结论交还现实：本节只建立功能层关系，不由此推出特定脑区实现、主观意识或宇宙目的。可测状态、干预后果和明确失败条件，才是从原文直觉走向可检验理论的桥。

\section{说服、协作、授权与委派}
\begin{narration}
现在把镜头移到“说服、协作、授权与委派”。我们只沿本章已经取得的概念向前一步：先确定它在闭环里的角色，再检查它的尺度与失败边界。
\end{narration}

本节把“说服、协作、授权与委派”落实为跨主体 CSO 流通过角色、协议、外部记忆和制度沉降形成的慢尺度协调结构。这个定义不是说“任何相似现象都算”，而是要求我们同时给出对象域、允许的干预以及观察窗口。对于主题“具身外部记忆与社会性外部记忆”而言，它承担的是固定功能拓扑中的持续主体中的一个可辨识环节。

把这一关系写成最小数学骨架，可得

\begin{equation}\mathcal O=(\{A_i\},\mathcal P,\mathcal M,\mathcal N).\end{equation}

式中的符号只承担结构角色，具体参数仍需由任务和身体数据辨识。我们首先检查可观测量，再检查控制量，最后检查比它更慢的约束。这样的顺序来自持续学习、工作空间、观测器—控制器与混杂系统，但本书对说服、协作、授权与委派的命名和组合仍是需要实验验证的建模提案。

这一小节的反例同样重要：把组织类比成单一大脑会虚构不存在的中央观察器。因此评价它不能只看一次任务结果，而要比较干预前后的轨迹、恢复时间、维护成本和风险。如果改变该机制而所有允许观察下的分布都不变，那么它至少在当前实验族中是冗余描述。

\begin{thought}
对“说服、协作、授权与委派”做一次消融：冻结它、删除它，或只允许它以相邻时间尺度交换残差。哪些现象立即消失，哪些只在长期出现？若答案无法区分三种干预，我们还没有找到正确的状态变量。
\end{thought}

最后，把结论交还现实：本节只建立功能层关系，不由此推出特定脑区实现、主观意识或宇宙目的。可测状态、干预后果和明确失败条件，才是从原文直觉走向可检验理论的桥。

\section{群体认知}
\begin{narration}
现在把镜头移到“群体认知”。我们只沿本章已经取得的概念向前一步：先确定它在闭环里的角色，再检查它的尺度与失败边界。
\end{narration}

本节把“群体认知”落实为跨主体 CSO 流通过角色、协议、外部记忆和制度沉降形成的慢尺度协调结构。这个定义不是说“任何相似现象都算”，而是要求我们同时给出对象域、允许的干预以及观察窗口。对于主题“具身外部记忆与社会性外部记忆”而言，它承担的是固定功能拓扑中的持续主体中的一个可辨识环节。

把这一关系写成最小数学骨架，可得

\begin{equation}\mathcal O=(\{A_i\},\mathcal P,\mathcal M,\mathcal N).\end{equation}

式中的符号只承担结构角色，具体参数仍需由任务和身体数据辨识。我们首先检查可观测量，再检查控制量，最后检查比它更慢的约束。这样的顺序来自持续学习、工作空间、观测器—控制器与混杂系统，但本书对群体认知的命名和组合仍是需要实验验证的建模提案。

这一小节的反例同样重要：把组织类比成单一大脑会虚构不存在的中央观察器。因此评价它不能只看一次任务结果，而要比较干预前后的轨迹、恢复时间、维护成本和风险。如果改变该机制而所有允许观察下的分布都不变，那么它至少在当前实验族中是冗余描述。

\begin{thought}
对“群体认知”做一次消融：冻结它、删除它，或只允许它以相邻时间尺度交换残差。哪些现象立即消失，哪些只在长期出现？若答案无法区分三种干预，我们还没有找到正确的状态变量。
\end{thought}

最后，把结论交还现实：本节只建立功能层关系，不由此推出特定脑区实现、主观意识或宇宙目的。可测状态、干预后果和明确失败条件，才是从原文直觉走向可检验理论的桥。

\section{组织制度作为超慢外部结构}
\begin{narration}
现在把镜头移到“组织制度作为超慢外部结构”。我们只沿本章已经取得的概念向前一步：先确定它在闭环里的角色，再检查它的尺度与失败边界。
\end{narration}

本节把“组织制度作为超慢外部结构”落实为跨主体 CSO 流通过角色、协议、外部记忆和制度沉降形成的慢尺度协调结构。这个定义不是说“任何相似现象都算”，而是要求我们同时给出对象域、允许的干预以及观察窗口。对于主题“具身外部记忆与社会性外部记忆”而言，它承担的是固定功能拓扑中的持续主体中的一个可辨识环节。

把这一关系写成最小数学骨架，可得

\begin{equation}\mathcal O=(\{A_i\},\mathcal P,\mathcal M,\mathcal N).\end{equation}

式中的符号只承担结构角色，具体参数仍需由任务和身体数据辨识。我们首先检查可观测量，再检查控制量，最后检查比它更慢的约束。这样的顺序来自持续学习、工作空间、观测器—控制器与混杂系统，但本书对组织制度作为超慢外部结构的命名和组合仍是需要实验验证的建模提案。

这一小节的反例同样重要：把组织类比成单一大脑会虚构不存在的中央观察器。因此评价它不能只看一次任务结果，而要比较干预前后的轨迹、恢复时间、维护成本和风险。如果改变该机制而所有允许观察下的分布都不变，那么它至少在当前实验族中是冗余描述。

\begin{thought}
对“组织制度作为超慢外部结构”做一次消融：冻结它、删除它，或只允许它以相邻时间尺度交换残差。哪些现象立即消失，哪些只在长期出现？若答案无法区分三种干预，我们还没有找到正确的状态变量。
\end{thought}

最后，把结论交还现实：本节只建立功能层关系，不由此推出特定脑区实现、主观意识或宇宙目的。可测状态、干预后果和明确失败条件，才是从原文直觉走向可检验理论的桥。

\section{个体—组织—文明的多尺度外化}
\begin{narration}
现在把镜头移到“个体—组织—文明的多尺度外化”。我们只沿本章已经取得的概念向前一步：先确定它在闭环里的角色，再检查它的尺度与失败边界。
\end{narration}

本节把“个体—组织—文明的多尺度外化”落实为承载范围、有效信息率与反馈周期之间的约束关系。这个定义不是说“任何相似现象都算”，而是要求我们同时给出对象域、允许的干预以及观察窗口。对于主题“具身外部记忆与社会性外部记忆”而言，它承担的是固定功能拓扑中的持续主体中的一个可辨识环节。

把这一关系写成最小数学骨架，可得

\begin{equation}T_{\min}\geq L/v+I/B+T_{\rm act}.\end{equation}

式中的符号只承担结构角色，具体参数仍需由任务和身体数据辨识。我们首先检查可观测量，再检查控制量，最后检查比它更慢的约束。这样的顺序来自持续学习、工作空间、观测器—控制器与混杂系统，但本书对个体—组织—文明的多尺度外化的命名和组合仍是需要实验验证的建模提案。

这一小节的反例同样重要：忽略传播和排队成本会制造虚假的实时性。因此评价它不能只看一次任务结果，而要比较干预前后的轨迹、恢复时间、维护成本和风险。如果改变该机制而所有允许观察下的分布都不变，那么它至少在当前实验族中是冗余描述。

\begin{thought}
对“个体—组织—文明的多尺度外化”做一次消融：冻结它、删除它，或只允许它以相邻时间尺度交换残差。哪些现象立即消失，哪些只在长期出现？若答案无法区分三种干预，我们还没有找到正确的状态变量。
\end{thought}

最后，把结论交还现实：本节只建立功能层关系，不由此推出特定脑区实现、主观意识或宇宙目的。可测状态、干预后果和明确失败条件，才是从原文直觉走向可检验理论的桥。

\section*{本章收束：把名词还原为闭环}
现在回头看“具身外部记忆与社会性外部记忆”，最重要的不是记住缩写，而是说清本章对象怎样改变系统的状态、代价或可恢复性。下一章“连续动力学与离散治理”只使用这里已经给出的结果，不反过来为本章提供前提。

\begin{bookproposition}
若一个候选机制既不改变可观测轨迹分布，也不改变资源、风险或恢复代价，那么在当前观察尺度上，它不是一个可辨识的功能结构。
\end{bookproposition}
\begin{proofidea}
功能等价由可干预后果定义。若所有允许干预下的输出分布与代价均不变，则该机制相对于当前实验族不可辨识；要么扩大观察尺度，要么删除这一冗余描述。
\end{proofidea}

\begin{readercheck}
请尝试回答：本章对象的状态是什么？它的最快和最慢时间常数是什么？什么观测能证伪它？它失败时，残差应向哪一层升级？
\end{readercheck}
\cleardoublepage
